%% 00_abstract.tex — Title, authorship, abstract, keywords.

\title{Efficiency of Asynchronous Dynamic Application Security Testing:\\
Payload Scheduling, Runtime Confirmation, and LLM-Based Triage\\
for False-Positive Reduction}

\author{%
\IEEEauthorblockN{Luis Ernesto Gald\'amez Vides}
\IEEEauthorblockA{%
Universidad Evang\'elica de El Salvador (UEES)\\
San Salvador, El Salvador\\
\texttt{2026010235@cvirtualuees.edu.sv}%
}
}

\maketitle

\begin{abstract}
Dynamic Application Security Testing (DAST) scanners are widely used in
continuous-integration pipelines, but their practical adoption is limited by
three factors: the accumulated latency of thousands of sequential requests,
the noise produced by a high false-positive rate, and the combinatorial
explosion of the payload space per injection point. This paper presents and
evaluates \herramienta, a DAST engine built entirely on \code{asyncio} that
introduces (i)~a payload-scheduling pipeline --- context interleaving,
a-priori confidence prioritization, and configurable ordering, including a
live adaptive-sorting feedback loop --- (ii)~a runtime-confirmation stage that
relabels every finding with a calibrated confidence level before it is
reported, and (iii)~an optional LLM-based triage stage that re-scores each raw
finding using request/response context, without an additional network round
trip to the target. The central hypothesis is that
these techniques improve the detection-to-cost ratio (number of requests)
without increasing the false-positive rate. To test it we built a reproducible
testbed on OWASP Benchmark restricted to \num{338}~\code{GET}-parameter
targets (\num{194}~vulnerable instances and \num{144}~known traps), an
ablation study under a randomized complete block design (RCBD) with four
request budgets and six pipeline configurations, and a battery of
non-parametric tests (Friedman, Nemenyi post-hoc, paired Wilcoxon with
Benjamini--Hochberg correction, and Cliff's $\delta$ as effect size).
Reproducibility is guaranteed through a per-run manifest that fixes the
SHA-256 hash of the payload corpus, the \code{git} commit, and the global seed
of the random-number generator. The full heuristic pipeline reaches
\ValorRecallTecho{} recall at a \ValorFPRBaseline{} false-positive rate and
already saturates at \num{20}~requests per parameter. Of the three
heuristic-pipeline optimizations, only confidence-based prioritization has a
practically consequential effect --- and it is bounded to tight budgets:
context interleaving and runtime confirmation shift the request order in a
statistically detectable way (Friedman, $p \ll 0.001$) but with negligible
effect sizes ($\lvert\delta\rvert < 0.02$). Layering LLM-based triage on top
of the baseline pipeline, evaluated separately across the same budget grid,
suppressed \emph{zero} findings at every budget under its default confidence
threshold, adding no measurable wall-clock overhead to the scan; we report
this null result openly and attribute it to the benchmark's unusually
unambiguous trigger signals rather than to the triage stage being validated
as safe or effective in general. Run against OWASP ZAP 2.17.0 on the
identical target list under the same oracle, \code{baseline} matches ZAP's
recall exactly (both detect the same \num{137}/\num{194} instances), at the
cost of a higher false-positive rate on the known traps (ZAP: \num{0}\,\%).
It reaches that recall \num{11}$\times$ faster than ZAP at its saturating
budget. We conclude that, on a benchmark with regular trigger signals, the
value of payload scheduling is one of \emph{anticipation} (higher AUC under
a tight budget), not of a higher detection ceiling. The LLM-triage stage
remains to be stress-tested on a noisier target.
\end{abstract}

\begin{IEEEkeywords}
web application security, DAST, dynamic analysis, asynchronous programming,
false positives, payload scheduling, LLM triage, reproducible evaluation,
OWASP Benchmark.
\end{IEEEkeywords}
